% Glossary for the document
% Optional, but you must delete the call in main.tex
% MIT License
%
% Copyright (c) 2026 HumberASV
% Copyright (c) 2026 Carson Fujita
%
% Permission is hereby granted, free of charge, to any person obtaining a copy
% of this software and associated documentation files (the "Software"), to deal
% in the Software without restriction, including without limitation the rights
% to use, copy, modify, merge, publish, distribute, sublicense, and/or sell
% copies of the Software, and to permit persons to whom the Software is
% furnished to do so, subject to the following conditions:
%
% The above copyright notice and this permission notice shall be included in all
% copies or substantial portions of the Software.
%
% THE SOFTWARE IS PROVIDED "AS IS", WITHOUT WARRANTY OF ANY KIND, EXPRESS OR
% IMPLIED, INCLUDING BUT NOT LIMITED TO THE WARRANTIES OF MERCHANTABILITY,
% FITNESS FOR A PARTICULAR PURPOSE AND NONINFRINGEMENT. IN NO EVENT SHALL THE
% AUTHORS OR COPYRIGHT HOLDERS BE LIABLE FOR ANY CLAIM, DAMAGES OR OTHER
% LIABILITY, WHETHER IN AN ACTION OF CONTRACT, TORT OR OTHERWISE, ARISING FROM,
% OUT OF OR IN CONNECTION WITH THE SOFTWARE OR THE USE OR OTHER DEALINGS IN THE
% SOFTWARE.

\AtBeginEnvironment{thebibliography}{\renewcommand{\autocite}{\textcite}} % Change the default citation command to textcite for better readability in the bibliography
% right now the default is set to numeric, but you can change it to authoryear or any other style supported by biblatex. 
% Adjust the citation command accordingly if you change the style. 
% For example, if you switch to authoryear, you might want to use \cite instead of \textcite in the bibliography.
% \AtBeginEnvironment{thebibliography}{\renewcommand{\autocite}{\cite}} % Uncomment this line if you prefer longer citations instead

% Note: Make sure you add the relevant entries to your .bib file for the citations used in the glossary.
\section{Glossary}
\begin{description}
	\item[\textbf{Autonomous Surface Vehicle (ASV)}]	
	An autonomous surface vehicle (ASV) is an unmanned vessel that operates on the sea surface without real-time input or control from human operators.\autocite{TREMBANIS2021251}

	\item[\textbf{Base Station}]	
	A fixed radio communications station that can communicate with mobile radio devices in its neighbourhood. Forms of radio communication that have base stations include radio wireless data and mobile telephony...\autocite{BASESTATION2018}
	
	\item[\textbf{DDS}]
	Short for Data Distribution Service a middleware for real-time comminication and security.\autocite{MATLABDDS}

	\item[\textbf{GPS}]
	Global Positioning System; "a navigational system using satellite signals to fix the location of a radio receiver on or above the earth's surface."\autocite{GPSDICTIONARY} 

	\item[\textbf{HumberASV}]	
	Humber Polytechnic's Mechatronics club. A multidisciplinary team of passionate Humber Polytechnic students and professionals dedicated to advancing autonomous maritime technology.\autocite{HUMBERASVTEAM}

	\item[\textbf{Hyperlink}]
	hyperlink, a link between related pieces of information by electronic connections in order to allow a user easy access between them.\autocite{HYPERLINKDICTIONARY}

	\item[\textbf{LaTex}]
	LaTeX, software used for typesetting technical documents.\autocite{LATEXDICTIONARY}
	
	\item[\textbf{LIDAR}]
	LiDAR, an acronym for "light detection and ranging," is a remote-sensing technology that uses laser beams to measure precise distances and movement in an environment, in real time.\autocite{LIDAR}

	\item[\textbf{Loon-E}]	
	Humber ASV's current ASV.\autocite{HUMBERASVTEAM}
	
	\item[\textbf{MQTT}]
	Message Queuing Telemetry Transport (MQTT) is a lightweight, publish-subscribe network protocol that transports messages between devices.\autocite{MQTT}

	\item[\textbf{Node}]
	A node is a process that performs computation \autocite{NODEDEFINITION}.

	\item[\textbf{Robot Operating System (ROS)}]	
	The Robot Operating System (ROS) is a set of software libraries and tools that help you build robot applications.\autocite{ROS}

	\item[\textbf{Pose}]
	In the fields of computing and computer vision, pose (or spatial pose) represents the position and the orientation of an object, each usually in three dimensions. \autocite{HoffNguyenLyon1996}

	\item[\textbf{PWM}]
	Acronym for Pulse Width Modulation.\autocite{PWMDICTIONARY}

	\item[\textbf{ROS2}]
	The Robot Operating System (ROS) is a set of software libraries and tools that help you build robot applications. From drivers to state-of-the-art algorithms, and with powerful developer tools, ROS has what you need for your next robotics project. And it's all open source \autocite{ROS2Humble}.
	
	\item[\textbf{RTP}]
	Real-time Transport Protocol.

	\item[\textbf{TCP}]
	Acronym for Transmission Control Protocol. TCP/IP, standard Internet communications protocols that allow digital computers to communicate over long distances. The Internet is a packet-switched network, in which information is broken down into small packets, sent individually over many different routes at the same time, and then reassembled at the receiving end. TCP is the component that collects and reassembles the packets of data, while IP is responsible for making sure the packets are sent to the right destination.\autocite{TCPIPDICTIONARY}

	\item[\textbf{TWIST}]
	Steering commands for a robot, consisting of linear and angular velocities.

	\item[\textbf{Topic}]
	Topics are named buses over which nodes exchange messages \autocite{TOPICDEFINITION}.

	\item[\textbf{Web Client}]
	See \hyperref[Web Browser]{Web Browser}.

	\item[\textbf{Web Browser}]
	Software that allows a computer user to find and view information on the Internet \autocite{BROWSERDICTIONARY}.

	\item[\textbf{Web Server}]
	server, network computer, computer program, or device that processes requests from a client \autocite{WEBSERVERDICTIONARY}.

\end{description}